% Page budget: exactly 1 page; no figures; no bullets; organize by method family.
\clearpage
\section{Related Work}

\subsection{Prediction-First Scheduling}
% [作者手写区] 要点: 归类持续信任预测的 length/rank 方法；每篇只写与本研究缺口直接相关的差异。Bib keys: fuLTR2024, ssjf2024, pars2025, egtp2026, elis2025.

\subsection{Prediction-Free and Preemptive Scheduling}
% [作者手写区] 要点: 归类不用输出预测或依靠运行时反馈/抢占的方法；显式区分模拟器零抢占成本与真实 KV 移动成本。Bib keys: fastServe2026, agentix2026, infercept2024.

\subsection{Corrected and Uncertainty-Aware Prediction}
% [作者手写区] 要点: 归类 tail-aware、uncertainty-aware、SLO-aware 修正；落点是“何时不该使用预测”的未闭合系统问题。Bib keys: tie2026, beyondPrediction2026, jitServe2026.

\subsection{Agentic Serving Systems}
% [作者手写区] 要点: 归类 workflow/KV/session 级系统，说明它们与本研究的 gateway-level reliability gate 是互补关系；不要声称复现未复现系统。Bib keys: agentix2026, infercept2024, bfcl, toolathlon2025.
